\subsection{The oriented smooth type}

\paragraph{Theorem F.4.1 (Proved: a complex structure on the stated connected sum).}
The compact complex fourfold of Theorem F.3.7 is three-connected,
has Euler number six, and is orientation-preservingly diffeomorphic to
\[
 3\mathbb H\mathbb P^2\ \#\ (-\mathbb H\mathbb P^2).
 \tag{F.25}
\]
Its integral middle lattice is $I_{3,1}$ and, in a diagonal integral
basis, $p_1/2=(1,1,1,1)$.  In particular the conclusion is an exact
smooth diffeomorphism, with no remaining homotopy-sphere ambiguity.

\begin{proof}
Theorem F.2.6 gives three-connectivity and free integral homology
with middle rank four.  The complex orientation and Theorem F.1.2
give signature two, the odd unimodular form $I_{3,1}$, and
$c_2^2=2$.  Since $H^2=0$, we have $c_1=0$ and
$p_1=-2c_2$.  The reduction of $c_2$ modulo two is $w_4$.
The Wu formula, with $w_1=w_2=0$, says that $w_4$ is characteristic
for the middle form.  Thus
$v=p_1/2=-c_2$ is characteristic and $v^2=2$.
It is primitive: a nontrivial integral divisor would have square
dividing two.  This excludes every divisor of absolute value
at least two.

We verify the orbit of this vector in $L=I_{3,1}$.
Its orthogonal complement $K=v^\perp$ is even, indefinite of
signature $(2,1)$, with determinant of absolute value two.
Unimodularity of $L$ and primitivity of $v$ give $w\in L$ with
$(w,v)=1$.  The vector $k=w-v/2$ represents the nonzero element
of $K^*/K\simeq\mathbb Z/2$.  Its discriminant square is
\[
 k^2=w^2-\frac12\equiv\frac12\pmod {2\mathbb Z},
 \tag{F.26}
\]
because characteristicity makes $w^2$ odd.
Consequently $K$ has the signature and discriminant form of
$U\oplus\langle2\rangle$, hence the same genus by
\cite[Corollary 1.9.4]{ffNikulin}.
Nikulin's Theorem 1.14.2 applies: the rank is three, the
odd-primary discriminant lengths are zero, and the two-primary
length is one, so its additional equality-case condition at two
does not arise.  The genus has a single isomorphism class.
Thus $K\simeq U\oplus\langle2\rangle$.
The extension from $\mathbb Zv\oplus K$ to the odd unimodular
lattice $L$ has index two.  Its gluing subgroup is the unique
nontrivial graph between the two discriminant groups of order
two.  Every isomorphism of complements therefore extends, taking
$v$ to a chosen characteristic vector of square two.
The vector $(1,1,1,1)$ is such a vector in $I_{3,1}$.
This proves the asserted orbit, including the class of $p_1$.

Let $M=3\mathbb H\mathbb P^2\#(-\mathbb H\mathbb P^2)$ with
the displayed orientation.  Its intersection form is $I_{3,1}$,
and its standard generators have $p_1=2\sum_i x_i$, as used
and derived in Theorem F.1.3.  Therefore the two manifolds have
isomorphic middle forms and first Pontryagin classes.
These determine Wall's full normal invariant in dimension eight.
Indeed the normal rank-four bundle of a representative embedded
$S^4$ has Euler number its self-intersection and Pontryagin
number the evaluation of $p_1(TM)$, since $p_1(TS^4)=0$.
The clutching group is
$\pi_3\operatorname{Spin}(4)=\mathbb Z\oplus\mathbb Z$;
with compatible generators $(a,b)$ its two characteristic
numbers are $e=a-b$ and $p_1=-2(a+b)$.
This pair is injective, so no additional unstable normal
invariant is omitted.

Wall's classification theorem in the precise formulation
\cite[Recollection 2.2, Definition 2.5 and Theorem 2.6]{ffBHS}
now gives an oriented diffeomorphism after connected sum with
some $\Sigma\in\Theta_8$.  We remove this ambiguity using the
exact inertia calculation
\cite[Theorem 1.4 and its proof in Section 9]{ffNagy2026}:
for a closed three-connected eight-manifold its inertia group
is zero exactly when $p_1$ is divisible by eight; otherwise,
since $\Theta_8=\mathbb Z/2$, it is all of $\Theta_8$.
Here $p_1$ has divisibility two.  Thus $M\#\Sigma$ is oriented
diffeomorphic to $M$ for both possibilities for $\Sigma$.
This proves (F.25).
\end{proof}

The application of Nagy's theorem uses the published 2026 paper
by Csaba Nagy, whose theorem and proof were opened in this run;
it is not an inference from the earlier announcement with
Crowley.  Mandal's almost-complex criterion alone gives no
integrability statement.  Here integrability comes from the
periods and the compact charts of Theorems F.3.5--F.3.7, while
the cited classification results identify their smooth manifold.
